% File: sections/00_abstract.tex

\noindent This paper proposes \textbf{Collective}, a decentralized architecture for social coordination designed to supersede the competitive, individualistic logic of traditional market economies. By conceptualizing the economy as a form of \textit{ubiquitous computation}—a distributed virtual machine for the satisfaction of community demand—we move beyond the constraints of private ownership and wage-based labor. 

The system is architected as a federated mesh of recursive \textbf{Service Nodes}, which utilize a Work Breakdown Structure (WBS) modeled via Business Process Model and Notation (BPMN) to map community needs to executable tasks. We introduce a dual-token metabolic model: \textit{Value}, a liquid credit for lifestyle construction that facilitates unbounded negative balances for social security; and \textit{Tokens}, a non-transferable reputation metric earned through labor and required for node-level governance. 

The structural integrity of the mesh is maintained by \textbf{Societies}, specialized nodes that represent the legal and territorial logic of the community, governing resource stewardship and member rights. Central to the model is the \textbf{Common Fund}, which mediates the "lease" of external resources to allow Collective cells to coexist with and eventually "leach" the legacy market. Through the integration of \textbf{Bodies of Knowledge} for skill verification and \textbf{Prospecting Engines} for demand forecasting, Collective provides a scalable framework for a post-scarcity metabolism. Ultimately, the system transforms the daily actions of human interaction into a self-organizing protocol, enabling the phased transition toward a society where contribution is the sole investment and utility is the primary output.